\subsection{两个测度的乘积}

设 S 和 T 是两个拓扑空间，分别赋有两个（正的）预测度 \(\lambda\) 和 \(\mu\)，令 X 为乘积空间 \(S \times T\)。设 K 是 X 的紧子集；以 A 和 B 分别表示 K 在 S 和 T 上的投影，并置

\[
\nu_ {\mathrm{K}} = \left(\lambda_ {\mathrm{A}} \otimes \mu_ {\mathrm{B}}\right) _ {\mathrm{K}}.\tag{3}
\]

由此在 X 上定义了一个预测度。事实上，令 L 是 X 中包含 K 的紧子集，令 C 和 D 是它的两个投影；则 \(A \subset C\)、\(B \subset D\)，因而利用诱导测度的传递性以及 Ch. V, §8, No.~5 的 Prop. 12，有

\[
\begin{array}{r l} & {(\nu_ {\mathrm{L}}) _ {\mathrm{K}} = \big ((\lambda_ {\mathrm{C}} \otimes \mu_ {\mathrm{D}}) _ {\mathrm{L}} \big) _ {\mathrm{K}} = (\lambda_ {\mathrm{C}} \otimes \mu_ {\mathrm{D}}) _ {\mathrm{K}}} \\ & {\qquad = \big ((\lambda_ {\mathrm{C}} \otimes \mu_ {\mathrm{D}}) _ {\mathrm{A} \times \mathrm{B}} \big) _ {\mathrm{K}} = (\lambda_ {\mathrm{A}} \otimes \mu_ {\mathrm{B}}) _ {\mathrm{K}} = \nu_ {\mathrm{K}}.} \end{array}
\]

定义 6. --- 由 (3) 定义的预测度 \(\nu\) 称为 \(\lambda\) 和 \(\mu\) 的乘积预测度，记作 \(\lambda \otimes \mu\)。

这个定义显然可以推广到 \(\lambda\) 和 \(\mu\) 是复预测度的情形；此时有 \(|\lambda\otimes\mu|=|\lambda|\otimes|\mu|\)（Ch. III, §4, No.~2, Prop. 3 and Ch. IV, §5, No.~7, Lemma 3）。

若 \(f\) 和 \(g\) 分别是定义在 \(\overline{\mathbf{R}}_+\) 和 \(\mathbf{C}\) 上的函数，则函数 \((s,t)\mapsto f(s)g(t)\) 定义在 \(S\times T\) 上，记作 \(f\otimes g\)。

命题 10. --- 设 \(\nu\) 是 \(\lambda\) 和 \(\mu\) 的乘积预测度；对于每个函数 \(f \in \mathcal{F}_{+}(\mathrm{S})\) 和每个函数 \(g \in \mathcal{F}_{+}(\mathrm{T})\)，

\[
\nu^ {\bullet} (f \otimes g) = \lambda^ {\bullet} (f) \mu^ {\bullet} (g).\tag{4}
\]

预测度 \(\nu\) 是满足 (4) 的 \(S \times T\) 上唯一的预测度。

当 K（相应地，L）遍历 S（相应地，T）的紧子集时，有

\[
\begin{array}{l} \nu^ {\bullet} (f \otimes g) = \sup _ {K, L} \nu_ {K \times L} ^ {\bullet} \big ((f \otimes g) _ {K \times L} \big) = \sup _ {K, L} (\lambda_ {K} \otimes \mu_ {L}) ^ {\bullet} (f _ {K} \otimes g _ {L}) \\ \qquad = \sup _ {K, L} \lambda_ {K} ^ {\bullet} (f _ {K}) \cdot \mu_ {L} ^ {\bullet} (g _ {L}) = \big (\sup _ {K} \lambda_ {K} ^ {\bullet} (f _ {K}) \big) \big (\sup _ {L} \mu_ {L} ^ {\bullet} (g _ {L}) \big) \\ \qquad = \lambda^ {\bullet} (f) \mu^ {\bullet} (g) \end{array}
\]

这是由 Ch. V, §8, No.~3 的 Prop. 8 得到的。

设 \(\eta\) 是 \(S \times T\) 上满足 (4) 的另一个预测度，设 \(K\) 和 \(L\) 分别是 \(S\) 和 \(T\) 的紧子集，\(f\) 和 \(g\) 分别是 \(\mathcal{F}_{+}(K)\) 和 \(\mathcal{F}_{+}(L)\) 中的元素。对于延拓为 0 的函数，有关系 \((f \otimes g)^{0} = f^{0} \otimes g^{0}\)，因此（§1, No.~2, Prop. 2）

\[
\begin{array}{r l} & {\eta_ {\mathrm{K} \times \mathrm{L}} ^ {\bullet} (f \otimes g) = \eta^ {\bullet} \big ((f \otimes g) ^ {0} \big) = \eta^ {\bullet} (f ^ {0} \otimes g ^ {0})} \\ & {\qquad = \lambda^ {\bullet} (f ^ {0}) \mu^ {\bullet} (g ^ {0}) = \lambda_ {\mathrm{K}} ^ {\bullet} (f) \mu_ {\mathrm{L}} ^ {\bullet} (g).} \end{array}
\]

特别地，若取 \(f \in \mathcal{K}_{+}(\mathrm{K})\)、\(g \in \mathcal{K}_{+}(\mathrm{L})\)，则可见 \(\eta_{\mathrm{K} \times \mathrm{L}}\) 具有乘积测度 \(\lambda_{\mathrm{K}} \otimes \mu_{\mathrm{L}}\) 的特征性质（Ch. III, §4, No.~1, Th. 1）。因此 \(\eta_{\mathrm{K} \times \mathrm{L}} = \nu_{\mathrm{K} \times \mathrm{L}}\)；由于 \(S \times T\) 的每个紧子集都包含于某个形如 \(K \times L\) 的集合中，诱导测度的传递性蕴含 \(\eta = \nu\)。

推论 1. --- 如果 \(\lambda\) 和 \(\mu\) 是测度，则 \(\nu\) 是测度。

事实上，令 \(x = (s, t)\) 是 X 中的一点，令 U 和 V 分别是 s 和 t 的邻域，并满足 \(\lambda^{\bullet}(\mathrm{U}) < +\infty\)、\(\mu^{\bullet}(\mathrm{V}) < +\infty\)；集合 \(U \times V\) 是 x 的邻域，并且由 (4) 有 \(\nu^{\bullet}(\mathrm{U} \times \mathrm{V}) = \lambda^{\bullet}(\mathrm{U}) \mu^{\bullet}(\mathrm{V}) < +\infty\)；因此负荷 \(\nu^{\bullet}\) 局部有界，预测度 \(\nu\) 是测度。

这个结果立即推广到复测度。

推论 2. --- 如果 A 是 S 的对 \(\lambda\) 局部可忽略的子集，则 \(A \times T\) 对 \(\nu\) 局部可忽略。

推论 3. --- 假设 \(\lambda\)（相应地，\(\mu\)）是 S（相应地，T）上测度的可求和族 \((\lambda_{\alpha})_{\alpha \in A}\)（相应地，\((\mu_{\beta})_{\beta \in B}\)）之和。则族 \((\lambda_{\alpha} \otimes \mu_{\beta})_{(\alpha, \beta) \in A \times B}\) 可求和，且其和为 \(\lambda \otimes \mu\)。

事实上，令 \(p\) 为负荷 \(\sum_{\alpha, \beta} (\lambda_{\alpha} \otimes \mu_{\beta})^{\bullet}\)；若 \(f \in \mathcal{F}_{+}(S)\) 且 \(g \in \mathcal{F}_{+}(T)\)，则显然有 \(p(f \otimes g) = \lambda^{\bullet}(f)\mu^{\bullet}(g)\)。推论 1 的证明表明 \(p\) 局部有界，因此族 \((\lambda_{\alpha} \otimes \mu_{\beta})\) 可求和（§1, No.~7, Prop. 7）。其和 \(\eta\) 满足 \(\eta^{\bullet} = p\)（§1, No.~7, Prop. 7），而 Prop. 10 蕴含 \(\eta = \nu\)。

